\documentclass[11pt, a4paper]{article}
\usepackage[utf8]{inputenc}
\usepackage{amsmath, amssymb}
\usepackage{graphicx}
\usepackage{hyperref}
\usepackage{xcolor}
\usepackage{listings}
\lstset{
    language=Python,
    basicstyle=\ttfamily\small,
    keywordstyle=\color{blue}\bfseries,
    commentstyle=\color{gray}\itshape,
    stringstyle=\color{orange},
    numberstyle=\tiny\color{gray},
    numbers=left,
    stepnumber=1,
    numbersep=8pt,
    frame=single,
    breaklines=true,
    breakatwhitespace=false,
    showstringspaces=false,
    tabsize=4,
    captionpos=b
}
\usepackage{geometry}
\geometry{margin=1in}

\title{EE 325 Digital Signal Processing \\ Assignment 1: Sampling and Aliasing}
\author{PERERA J.D.T. \\ E/21/291}
\date{\today}

\begin{document}
\maketitle

\section{Objective}
Demonstrate the effects of sampling a continuous-time signal at frequencies below the Nyquist rate (Aliasing).

\section{Problem Description}
\textbf{Continuous Signal:}
A sine wave of frequency $f = 1000$ Hz is used:
$$x_c(t) = \sin(2\pi \cdot 1000 t)$$
\textbf{Nyquist Rate:}
The minimum sampling rate required to avoid aliasing is $2f_{max} = 2 \times 1000 = 2000$ Hz.

\section{Simulation Analysis}
The Python script sampled this signal at two different frequencies:

\subsection{Case 1: Sampling at $f_s = 600$ Hz}
\begin{itemize}
    \item \textbf{Observation:} Since $f_s = 600$ Hz $< 2000$ Hz, the system is \textbf{under-sampled}.
    \item \textbf{Result:} Aliasing occurs. The high-frequency original signal ($1000$ Hz) will appear as a lower frequency alias in the discrete domain.
    \item \textbf{Alias Frequency:} $|1000 - k \cdot 600|$. For $k=2$, $|1000 - 1200| = 200$ Hz (folded).
\end{itemize}

\subsection{Case 2: Sampling at $f_s = 1500$ Hz}
\begin{itemize}
    \item \textbf{Observation:} Since $f_s = 1500$ Hz $< 2000$ Hz, the system is also \textbf{under-sampled}.
    \item \textbf{Result:} Aliasing occurs.
    \item \textbf{Alias Frequency:} $|1000 - 1500| = 500$ Hz.
\end{itemize}

\section{Conclusion}
The simulation visually demonstrates that sampling at a rate lower than the Nyquist rate results in a distorted or ``aliased'' signal that does not faithfully represent the original frequency. To correctly reconstruct the 1000 Hz sine wave, a sampling rate of at least 2000 Hz is required.

\appendix
\section{Python Source Code}
\textbf{File:} \texttt{E21291\_Aliasing\_demo.py}
\begin{lstlisting}[caption={Aliasing Demonstration -- Assignment 1}]
#Aliasing Demonstration

import numpy as np
import matplotlib.pyplot as plt

# genarate continuous signal 
t = np.linspace(0, 0.01, 1000)  # time vector from 0 to 1 seconds
x_c = np.sin(2 * np.pi * 1000*t)  # 1000 Hz sine wave

# Sampling frequencies

f_s1 = 600
f_s2 = 1500

# Sampled signals

t_s1 = np.arange(0, 0.01, 1/f_s1)
x_s1 = np.sin(2 * np.pi * 1000*t_s1)

t_s2 = np.arange(0, 0.01, 1/f_s2)
x_s2 = np.sin(2 * np.pi * 1000*t_s2)

# Plotting

plt.figure(figsize=(12, 8))
plt.subplot(2, 1, 1)
plt.plot(t, x_c, label='Continuous Signal (1000 Hz)', color='blue')
plt.stem(t_s1, x_s1, linefmt='r-', markerfmt='ro', basefmt=' ', label='Sampled Signal (600 Hz)')
plt.title('Sampling at 600 Hz')
plt.xlabel('Time (s)')
plt.ylabel('Amplitude')
plt.legend()
plt.grid()

plt.subplot(2, 1, 2)
plt.plot(t, x_c, label='Continuous Signal (1000 Hz)', color='blue')
plt.stem(t_s2, x_s2, linefmt='g-', markerfmt='go', basefmt=' ', label='Sampled Signal (1500 Hz)')
plt.title('Sampling at 1500 Hz')
plt.xlabel('Time (s)')
plt.ylabel('Amplitude')
plt.legend()
plt.grid()

plt.tight_layout()
plt.show()

# discrete time waveform

n1 = np.arange(0, len(x_s1))
n2 = np.arange(0, len(x_s2))

plt.figure(figsize=(12, 6))
plt.subplot(2, 1, 1)
plt.stem(n1, x_s1, linefmt='r-', markerfmt='ro', basefmt=' ', label='Sampled Signal (600 Hz)')
plt.title('Discrete Time Signal at 600 Hz Sampling')
plt.xlabel('Sample Number')
plt.ylabel('Amplitude')
plt.legend()
plt.grid()

plt.subplot(2, 1, 2)
plt.stem(n2, x_s2, linefmt='g-', markerfmt='go', basefmt=' ', label='Sampled Signal (1500 Hz)')
plt.title('Discrete Time Signal at 1500 Hz Sampling')
plt.xlabel('Sample Number')
plt.ylabel('Amplitude')
plt.legend()
plt.grid()

plt.tight_layout()
plt.show()  

# comparsion of the sampled signals

plt.figure(figsize=(12, 6))
plt.plot(t, x_c, label='Continuous Signal (1000 Hz)', color='blue')
plt.stem(t_s1, x_s1, linefmt='r-', markerfmt='ro', basefmt=' ', label='Sampled Signal (600 Hz)')
plt.stem(t_s2, x_s2, linefmt='g-', markerfmt='go', basefmt=' ', label='Sampled Signal (1500 Hz)')
plt.title('Comparison of Sampled Signals')
plt.xlabel('Time (s)')
plt.ylabel('Amplitude')
plt.legend()
plt.grid()
plt.show()

print("The Nyquist rate for a 1000Hz signal is 2000Hz.")
print("When the 1000Hz signal is smapled below the Nyquist rate 600Hz, 1500Hz, aliasing occurs, distorting the signal representation.")
\end{lstlisting}

\end{document}
